\title{Controlling Text-to-Image Diffusion by Orthogonal Finetuning}

\begin{abstract}
State-of-the-art text-to-image diffusion models exhibit remarkable proficiency in producing lifelike images conditioned on textual descriptions. However, the challenge of systematically steering these advanced models to accommodate diverse downstream applications remains largely unresolved. In response, we present Orthogonal Finetuning~(OFT), a theoretically grounded adaptation strategy tailored for text-to-image diffusion models facing downstream tasks. Distinct from prior approaches, OFT is capable of rigorously maintaining hyperspherical energy, which encapsulates the inter-neuronal relationships when projected onto a unit hypersphere. This maintenance proves vital for retaining the semantic generative strength of text-to-image diffusion models. To further enhance the robustness of finetuning, we introduce Constrained Orthogonal Finetuning (COFT), an extension that enforces a fixed-radius constraint on the hypersphere. In this work, we concentrate on two critical adaptation scenarios for text-to-image finetuning: subject-driven generation—where a handful of exemplar images and a textual cue are used to generate new images of a specific subject—and controllable generation, which empowers the model with the capacity to utilize supplementary control signals as input. Empirical studies demonstrate that our OFT methodology consistently yields superior image generation quality and faster convergence compared to existing alternatives.
\end{abstract}

\section{Introduction}

Advances in text-to-image diffusion models~\cite{saharia2022photorealistic,ramesh2022hierarchical,rombach2022high} have enabled them to achieve exceptional performance in high-resolution image synthesis guided by textual inputs. Nevertheless, text-based conditioning often lacks precision and granularity, resulting in challenges for accurately guiding the resulting images. This work seeks to address this issue by focusing on two prominent tasks in text-to-image generation:

\begin{itemize}[leftmargin=*,nosep]

    \item \textbf{Subject-driven generation}~\cite{ruiz2023dreambooth}: Provided a limited set of images of a particular subject, this task aims to produce additional images depicting the same subject in varied scenarios, directed by a textual prompt.
    \item \textbf{Controllable generation}~\cite{zhang2023adding,mou2023t2i}: By incorporating supplementary control inputs (\emph{e.g.}, canny edge detections, segmentation masks), the objective is to synthesize images that respect both the control signal and an associated textual prompt.
\end{itemize}

Both of these tasks ultimately address the challenge of how to adapt text-to-image diffusion models through finetuning while maintaining the generative capabilities learned during pretraining. The requirements for an effective finetuning approach can be outlined as follows: (1) \emph{training efficiency}, characterized by a reduced number of trainable parameters and minimized training epochs, and (2) \emph{generalizability preservation}, referring to the retention of high-quality and diverse image generation. In practice, finetuning is generally performed either by slightly adjusting the model’s existing neuron weights using a small learning rate (\emph{e.g.}, \cite{ruiz2023dreambooth}), or by appending lightweight modules with reparameterized neuron weights to the model (\emph{e.g.}, \cite{hulora2022,zhang2023adding}). However, these straightforward strategies do not offer guarantees regarding the retention of pretraining generative performance. Moreover, the field lacks a theoretically grounded framework for crafting effective finetuning methods or for determining critical hyperparameters, such as the number of training epochs. One central obstacle is the absence of a metric that explicitly evaluates how well the pretrained model’s generative abilities are preserved after finetuning. Current approaches implicitly posit that reducing the Euclidean distance between the weights of the finetuned and pretrained models correlates with better preservation of pretrained capabilities. Consequently, extremely low learning rates are commonly employed. While this supposition occasionally holds true, we contend that relying solely on Euclidean distance does not provide a comprehensive measure of semantic retention. Therefore, introducing a more structural metric for assessing the divergence between the finetuned and pretrained models could substantially improve both the retention of generative performance and the overall stability of the finetuning process.

Building on prior findings that hyperspherical similarity effectively encapsulates semantic relationships~\cite{liu2017deep,liu2018decoupled,chen2020angular}, we leverage hyperspherical energy~\cite{liu2018learning} as a metric to quantify the pairwise interactions among neurons within a layer. Hyperspherical energy, calculated by aggregating hyperspherical similarities (\emph{e.g.}, cosine similarity) for all neuron pairs in a layer, reflects the extent of neuron distribution uniformity on the unit hypersphere~\cite{liu2021learning}. We posit that an optimal finetuned model should exhibit a hyperspherical energy level that closely aligns with its pretrained counterpart. While one straightforward approach involves adding a regularization term to maintain consistent hyperspherical energy throughout finetuning, this method does not ensure that deviations are strictly minimized. To address this, we exploit the invariance of hyperspherical energy under orthogonal transformations—specifically, if every neuron in a layer undergoes the same orthogonal transformation, the pairwise hyperspherical similarities remain invariant. Drawing on this property, we introduce Orthogonal Finetuning~(OFT), a technique designed to adapt large text-to-image diffusion models for downstream tasks without altering their hyperspherical energy. The fundamental principle of OFT is to train a single orthogonal transformation per layer that maintains the angular relationships among neurons. This can also be interpreted as recalibrating the neurons’ basis within each layer. Moreover, acknowledging that preserving the pretrained model’s performance is facilitated by minimizing Euclidean distance between the finetuned and pretrained weights, we extend our approach to Constrained Orthogonal Finetuning~(COFT), wherein the finetuned parameters are restricted to a fixed-radius hypersphere centered at the pretrained configuration.

The rationale behind utilizing orthogonal transformations for neuron finetuning is partially inspired by the properties of the 2D Fourier transform, which decomposes an image into its magnitude and phase spectra. The phase spectrum, capturing the angular relationship between the input and the basis, is largely responsible for retaining semantic content. Notably, reconstructing an image using its original phase spectrum alongside an arbitrary magnitude spectrum tends to preserve the core semantics. This observation implies that modifying neuron orientations—akin to adjusting the phase—plays a pivotal role in altering image semantics, which is central to both subject-driven and controllable generation tasks. However, allowing unrestricted modifications of neuron directions could compromise the generative capacity acquired during pretraining. To address this, we advocate for constraining the extent of such modifications by maintaining the angular relationships between all pairs of neurons, premised on the hypothesis that these angles encode essential neural network knowledge. Consequently, we introduce a strategy of learning a layer-shared orthogonal transformation for neurons within each layer, thereby preserving the hyperspherical energy throughout the process.

Additionally, our work is influenced by methods such as orthogonal over-parameterized training~\cite{liu2021orthogonal}, where classification neural networks are trained from scratch by applying orthogonal transformations to randomly initialized architectures. Such random initializations are expected to exhibit low hyperspherical energy, and the method of \cite{liu2021orthogonal} explicitly maintains this low energy to promote generalization in classification tasks~\cite{liu2018learning,lin2020regularizing}. Their findings indicate that orthogonal transformations provide sufficient adaptability to develop generalizable neural networks. In contrast, our focus shifts to the finetuning of text-to-image diffusion models, with the goal of enhancing controllability and optimizing generative performance for downstream tasks. There are two principal distinctions between OFT and the approach described in \cite{liu2021orthogonal}. Firstly, whereas \cite{liu2021orthogonal} aims to reduce hyperspherical energy, our OFT framework is designed to preserve the energy level present in the pretrained model, thereby safeguarding its underlying structure during finetuning—since reducing hyperspherical energy during diffusion model finetuning could disrupt inherent semantic representations. Secondly, OFT introduces mechanisms to minimize the divergence from the pretrained model, culminating in a constrained adaptation scheme, in contrast to the unconstrained paradigm of \cite{liu2021orthogonal}. The challenge in finetuning is to navigate the balance between flexibility and retention of pretrained knowledge, and we contend that OFT strikes this balance effectively. The main contributions of our work are as follows:

\begin{itemize}[leftmargin=*,nosep]

    \item We introduce Orthogonal Finetuning, a new approach for adapting text-to-image diffusion models to enhance their controllability. To ensure stable adaptation, we also present a constrained version that bounds the angular shift from the original pretrained model.
    \item Unlike previous finetuning strategies, OFT adjusts model weights while rigorously maintaining hyperspherical energy—a quantity we empirically identify as crucial for retaining the semantic generative capabilities of the pretrained network.
    \item We evaluate OFT on subject-driven and controllable generation scenarios, conducting extensive experiments that reveal marked gains over state-of-the-art methods in terms of output quality, training convergence, and robustness during finetuning. Furthermore, OFT demonstrates greater efficiency, converging effectively with far fewer images and training epochs.
    \item For the controllable generation setting, we develop a novel control consistency metric to quantify how well control signals are preserved. This metric involves extracting control cues from generated samples and comparing them to their original versions, further substantiating the robust controllability delivered by OFT.
\end{itemize}

\section{Related Work}

\textbf{Text-to-image diffusion models}. The landscape of text-to-image synthesis has advanced considerably~\cite{saharia2022photorealistic,ramesh2022hierarchical,rombach2022high,gu2022vector,nichol2022glide}, propelled by progress in diffusion-based generative architectures~\cite{ho2020denoising,song2021score,song2021denoising,dhariwal2021diffusion} and advances in vision-language representation learning~\cite{su2020vlbert,lu2019vilbert,radford2021learning,wang2021simvlm,li2022blip,li2023blip,alayrac2022flamingo,schuhmann2022laion}. Both GLIDE~\cite{nichol2022glide} and Imagen~\cite{saharia2022photorealistic} employ pixel-space diffusion models: GLIDE jointly trains a text encoder together with the diffusion model using paired data, whereas Imagen relies on a fixed, pretrained text encoder. Diffusion models in the latent space are exemplified by Stable Diffusion~\cite{rombach2022high} and DALL-E2~\cite{ramesh2022hierarchical}: Stable Diffusion leverages VQ-GAN~\cite{esser2021taming} for its visual codebook-based latent space, and DALL-E2 employs CLIP~\cite{radford2021learning} for a unified image-text embedding space. Outside the diffusion paradigm, generative adversarial networks~\cite{reed2016generative,zhang2017stackgan,xu2018attngan,li2019controllable} and autoregressive models~\cite{ramesh2021zero,ding2021cogview,wu2022nuwa,yuscaling} have also been explored for text-to-image generation. OFT is a model-agnostic finetuning framework, making it broadly compatible with any diffusion-based text-to-image model.

\textbf{Subject-driven generation}. To minimize changes to the subject during image editing, methods such as~\cite{avrahami2022blended,nichol2022glide} condition on user-provided masks. Inversion-based strategies~\cite{choi2021ilvr,dhariwal2021diffusion,ramesh2022hierarchical,gal2022image} can manipulate backgrounds or context while keeping the subject intact. Local and global modifications can also be achieved without input masks, as demonstrated by~\cite{hertz2022prompt}. However, these techniques often struggle to maintain subject-specific visual details. Pivotal Tuning~\cite{roich2022pivotal} adapts a generator near an inverted latent code with identity-preserving regularization, and~\cite{nitzan2022mystyle} creates a personalized face prior from individual image sets. Instance-level generation~\cite{casanova2021instance} can yield various depictions of a target, though sometimes at the cost of losing crucial instance details. DreamBooth~\cite{ruiz2023dreambooth} uses customized tokens and a handful of subject images to finetune diffusion models, enforcing both image reconstruction and class-level prior preservation. In our work, OFT adopts DreamBooth’s methodological setup but, crucially, replaces standard finetuning with orthogonal transformations rather than simple gradient-based updates with small learning rates.

\textbf{Controllable generation}. Image-to-image translation is fundamentally a type of controllable generation problem. Traditional approaches have relied heavily on conditional generative adversarial networks~\cite{isola2017image,zhu2017toward,wang2018high,park2019semantic,choi2018stargan} to address this task. More recently, diffusion models have been explored as alternative frameworks for image-to-image translation~\cite{saharia2022palette,voynov2022sketch,wang2022pretraining}. ControlNet~\cite{zhang2023adding} represents a significant advance by enabling fine-grained control of pretrained diffusion models through targeted finetuning and adaptation to new control signals, which results in highly effective controllable generation. In parallel, T2I-Adapter~\cite{mou2023t2i} similarly enhances control over image generation by finetuning diffusion models for stronger responsiveness to conditioning inputs. Building upon the experimental settings in \cite{zhang2023adding,mou2023t2i}, we introduce OFT as a method for finetuning pretrained diffusion models. Our approach consistently produces improved controllability even with reduced data requirements and a smaller number of trainable parameters. Notably, OFT does not incur any additional inference-time computation.

\textbf{Model finetuning}. Tailoring large-scale pretrained models for downstream applications through finetuning has gained substantial traction~\cite{devlin2018bert,brown2020language,he2022masked}. Adaptation strategies—such as those presented in~\cite{hulora2022,houlsby2019parameter,pfeiffer2020adapterfusion}—are an established area of study within natural language processing. Among these, LoRA~\cite{hulora2022} is particularly pertinent to OFT, as it imposes a low-rank constraint on weight updates during the finetuning stage. However, OFT differs fundamentally by utilizing a layer-shared orthogonal transformation that updates neuron weights multiplicatively, which theoretically preserves the angular relationships among neurons in a layer and leads to enhanced stability.

\section{Orthogonal Finetuning}

\subsection{Why Does Orthogonal Transformation Make Sense?}

We begin by examining the rationale for employing orthogonal transformations when finetuning text-to-image diffusion models. Our analysis addresses two key aspects: (1) the motivation for adjusting the angular direction of neurons, and (2) the justification for specifically using orthogonal transformations to accomplish these angle modifications.

Addressing the first question, we take cues from the empirical findings presented in \cite{liu2018decoupled,chen2020angular}, which suggest that angular differences in features are a robust indicator of semantic disparity. Prior work by SphereNet~\cite{liu2017deep} has demonstrated that when a neural network is constructed such that all neurons are normalized onto a unit hypersphere, the model exhibits both comparable representational capacity and potentially improved generalization. This observation suggests that the orientation of neuron vectors sufficiently encodes the essential information within the data. 

To underscore the significance of neuron angles, we present a simplified experiment depicted in Figure~\ref{fig:toy_ae}, where a basic convolutional autoencoder is trained on images of flowers. During the training phase, feature maps are produced using the standard inner product, with $z$ representing the output corresponding to the convolutional kernel $\bm{w}$ and sliding window input $\bm{x}$. In the evaluation phase, we assess three distinct methods for constructing the feature map: (a) the inner product applied in training, (b) solely the magnitude component, and (c) solely the angular (directional) component. Results shown in Figure~\ref{fig:toy_ae} illustrate that neuron direction alone nearly reconstructs the input images perfectly, whereas magnitude information fails to retain meaningful content. Notably, the cosine-based activation in (c) is not employed during the training process, which relies exclusively on the inner product. This outcome indicates that neuron direction is primarily responsible for encapsulating the semantic content of the input images. Consequently, adjusting neuron directions is likely the most effective approach for modifying image semantics.

Turning to the second question, an intuitive technique for altering neuron directions is to apply a uniform rotation or reflection across all neurons within a layer, leading naturally to the use of orthogonal transformations. While it is possible to consider more general angular modifications—where neurons undergo different rotational adjustments—our investigation reveals that orthogonal transformations strike an optimal balance between adaptability and structure. Furthermore, evidence in \cite{liu2021orthogonal} supports the sufficiency of orthogonal transformations for learning in neural networks. To further substantiate this point, we conduct a regularization experiment, leveraging the ControlNet~\cite{zhang2023adding} experimental setup for controllable image generation; outcomes are depicted in Figure~\ref{fig:ortho}. The results indicate that our standard OFT, which includes the orthogonality constraint, consistently maintains stability and achieves precise control by epoch 20. Conversely, OFT implemented without orthogonality fails both to generate plausible images and to exert effective control. This empirical finding underlines the crucial role played by the orthogonality constraint within OFT.

\subsection{General Framework}

Traditionally, the finetuning process updates either the entire model or specific layers via gradient descent with a modest learning rate, which implicitly limits how much the finetuned parameters can deviate from their pretrained counterparts. Essentially, standard finetuning attempts to adapt the model while minimizing $\thickmuskip=2mu \medmuskip=2mu \| \bm{M} - \bm{M}^0\|$, where $\bm{M}$ and $\bm{M}^0$ denote the weights after and before finetuning, respectively. While this soft regularization intuitively anchors the parameters to the pretrained solution, it may still permit more change than is desirable when finetuning large models. LoRA addresses this by explicitly constraining the rank of the weight updates, enforcing $\thickmuskip=2mu \medmuskip=2mu \text{rank}(\bm{M} - \bm{M}^0)=r'$, with $r'$ chosen to be small. In contrast, OFT imposes a constraint based on pairwise neuron similarity, specifically $\thickmuskip=2mu \medmuskip=2mu \| \text{HE}(\bm{M})-\text{HE}(\bm{M}^0) \|=0$, where the hyperspherical energy of the weight matrix is denoted $\text{HE}(\cdot)$. 

Consider a fully connected layer represented by $\thickmuskip=2mu \medmuskip=2mu \bm{W}=\{\bm{w}_1,\cdots,\bm{w}_n\}\in\mathbb{R}^{d\times n}$, where each neuron is $\bm{w}_i\in\mathbb{R}^{d}$ ($\bm{W}^0$ denotes pretrained weights). The layer outputs $\thickmuskip=2mu \medmuskip=2mu \bm{z}\in\mathbb{R}^n$ as $\thickmuskip=2mu \medmuskip=2mu \bm{z}=\bm{W}^\top\bm{x}$, where input $\thickmuskip=2mu \medmuskip=2mu \bm{x}\in\mathbb{R}^d$. In the context of OFT, the objective is to minimize the difference in hyperspherical energy between the finetuned and original parameters:

\begin{equation}\label{eq:oft_principle}
\footnotesize
\min_{\bm{W}} \left\| \text{HE}(\bm{W})-\text{HE}(\bm{W}^0)\right\|~~\Leftrightarrow~~\min_{\bm{W}} \bigg{\|} \sum_{i\neq j} \|\hat{\bm{w}}_i-\hat{\bm{w}}_j\|^{-1}-\sum_{i\neq j} \|\hat{\bm{w}}^0_i-\hat{\bm{w}}^0_j\|^{-1}\bigg{\|}
\end{equation}

Here, $\thickmuskip=2mu \medmuskip=2mu \hat{\bm{w}}_i:=\bm{w}_i/\|\bm{w}_i\|$ is the normalized $i$-th neuron, and the hyperspherical energy for the fully connected layer $\bm{W}$ is defined as $\thickmuskip=2mu \medmuskip=2mu \text{HE}(\bm{W}):= \sum_{i\neq j} \|\hat{\bm{w}}_i-\hat{\bm{w}}_j\|^{-1}$. Notably, the optimal (minimal) value of Eq.~\eqref{eq:oft_principle} is zero. This minimal value is achieved when $\bm{W}$ and $\bm{W}^0$ differ only by an orthogonal transformation—either a rotation or a reflection—namely, $\thickmuskip=2mu \medmuskip=2mu \bm{W}=\bm{R}\bm{W}^0$, where $\thickmuskip=2mu \medmuskip=2mu \bm{R}\in\mathbb{R}^{d\times d}$ is an orthogonal matrix (its determinant, being $1$ or $-1$, distinguishes rotation from reflection). This is the central principle of OFT: finetuning only needs to optimize orthogonal transformations—shared at the layer level—to rotate or reflect the neurons in

\begin{equation}\label{eq:oft_general}
    \footnotesize
    \bm{z}=\bm{W}^\top\bm{x}=(\bm{R}\cdot\bm{W}^0)^\top\bm{x},~~~\text{s.t.}~\bm{R}^\top\bm{R}=\bm{R}\bm{R}^\top=\bm{I}
\end{equation}

Here, $\bm{W}$ indicates the OFT-finetuned weight matrix, and $\bm{I}$ represents the identity matrix. The structure of OFT is depicted in Figure~\ref{fig:block}. In a manner analogous to LoRA's zero initialization, it is necessary that OFT commences finetuning directly from the pretrained parameters $\bm{W}^0$. To guarantee this behavior, the orthogonal matrix $\bm{R}$ is initially set as the identity matrix, ensuring that the starting point for finetuning coincides with the original pretrained weights. Ensuring $\bm{R}$ remains orthogonal can be accomplished using differentiable orthogonalization techniques, as elaborated in \cite{liu2021orthogonal,lezcano2019cheap}. We will subsequently elaborate on efficient computational strategies to maintain orthogonality.

\subsection{Efficient Orthogonal Parameterization}\label{sect:eff_ortho}

Typical orthogonalization approaches, such as the Gram-Schmidt process, although differentiable, generally entail high computational overhead in practice~\cite{liu2021orthogonal}. To enhance efficiency, this work employs Cayley parameterization for generating orthogonal matrices. Concretely, the orthogonal matrix is parameterized as $\thickmuskip=2mu \medmuskip=2mu \bm{R}=(\bm{I}+\bm{Q})(I-\bm{Q})^{-1}$, where $\bm{Q}$ is a skew-symmetric matrix such that $\thickmuskip=2mu \medmuskip=2mu \bm{Q}=-\bm{Q}^\top$. The Cayley transform is more computationally tractable but introduces a mild restriction—it is only capable of yielding orthogonal matrices with determinant $1$, i.e., matrices contained in the special orthogonal group. Our experiments demonstrate, however, that this constraint does not negatively impact practical outcomes. Even with the Cayley transformation, the parameter count for $\bm{R}$ may still become excessive for large $d$. To improve parameter efficiency, we advocate representing $\bm{R}$ as a block-diagonal matrix comprising $r$ blocks, resulting in the following structure:

\begin{equation}
    \footnotesize
    \bm{R}=\text{diag}(\bm{R}_1,\bm{R}_2,\cdots,\bm{R}_r)=\begin{bmatrix}
    \bm{R}_1\in \text{O}(\frac{d}{r}) &  &\\
    &\ddots & \\
    & & \bm{R}_r\in \text{O}(\frac{d}{r})
    \end{bmatrix}\in \text{O}(d)
\end{equation}

Here, $\text{O}(d)$ represents the orthogonal group in $d$ dimensions, with $\thickmuskip=2mu \medmuskip=2mu \bm{R}\in\mathbb{R}^{d\times d}$ and each $\thickmuskip=2mu \medmuskip=2mu \bm{R}_i\in\mathbb{R}^{d/r\times d/r}$ for all $i$ signifying orthogonal matrices. In the specific case where $\thickmuskip=2mu \medmuskip=2mu r=1$, the block-diagonal orthogonal matrix reduces to an ordinary, fully unconstrained orthogonal matrix. The number of free parameters in a $d\times d$ orthogonal matrix equals $\thickmuskip=2mu \medmuskip=2mu d(d-1)/2$, yielding computational complexity of $\mathcal{O}(d^2)$. For an orthogonal matrix with $r$ blocks along the diagonal, the total parameter count becomes $\thickmuskip=2mu \medmuskip=2mu d(d/r-1)/2$, resulting in complexity $\thickmuskip=2mu \medmuskip=2mu \mathcal{O}(d^2/r)$. To further enhance parameter efficiency, one may enforce block sharing by setting $\thickmuskip=2mu \medmuskip=2mu \bm{R}_i=\bm{R}_j$ for any $i\neq j$, thereby decreasing parameter complexity to $\mathcal{O}(d^2/r^2)$. All these approaches preserve the orthogonality of the matrix $\bm{R}$, thus maintaining the property of conserving hyperspherical energy.

We proceed to compare the parameter efficiency of OFT and LoRA. LoRA, when parameterized with low rank $r'$, has $r'(d+n)$ parameters that require training. Should both $r$ and $r'$ be functions of $d$ (such as $r=r'=\alpha d$ with $\thickmuskip=2mu \medmuskip=2mu 0<\alpha\leq 1$), then the parameter complexity for LoRA increases to $\mathcal{O}(d^2+dn)$, whereas OFT achieves $\mathcal{O}(d)$. A practical comparison highlights these distinctions: for a weight matrix sized $\thickmuskip=2mu \medmuskip=2mu 128\times 128$, LoRA (using $r'=8$) contains $2,048$ trainable parameters, whereas OFT (with $r=8$ and without block sharing) reduces this count to $960$.

\subsection{Constrained Orthogonal Finetuning}

Original OFT can be further regularized by confining the finetuned parameters to remain within a specified distance from the pretrained parameters. COFT introduces this restriction by imposing the following transformation in the forward computation:

\begin{equation}\label{eq:coft_general}
    \footnotesize
    \bm{z}=\bm{W}^\top\bm{x}=(\bm{R}\cdot\bm{W}^0)^\top\bm{x},~~~\text{s.t.}~\bm{R}^\top\bm{R}=\bm{R}\bm{R}^\top=\bm{I},~\norm{\bm{R}-\bm{I}}\leq \epsilon
\end{equation}

This constraint enforces both orthogonality and a maximum $\epsilon$-proximity to the identity matrix. Cayley parameterization, reviewed in Section~\ref{sect:eff_ortho}, enables the orthogonality constraint. However, it is not straightforward to directly impose the $\epsilon$-proximity condition on matrices parameterized via the Cayley transform. To explore this, one can use a Neumann series expansion to approximate $\thickmuskip=2mu \medmuskip=2mu \bm{R}=(\bm{I}+\bm{Q})(I-\bm{Q})^{-1}$ as $\thickmuskip=2mu \medmuskip=2mu \bm{R}\approx\bm{I}+2\bm{Q}+\mathcal{O}(\bm{Q}^2)$, assuming the validity of the Neumann expansion.

As a result, the requirement $\thickmuskip=2mu \medmuskip=2mu\|\bm{R}-\bm{I}\|\leq\epsilon$ can be equivalently reformulated within the Cayley transform, leading to the alternative condition $\thickmuskip=2mu \medmuskip=2mu \|\bm{Q}-\bm{0}\|\leq\epsilon'$ where $\bm{0}$ represents the zero matrix and $\epsilon'$ is a distinct hyperparameter controlling the allowable deviation. This new constraint imposed on the matrix $\bm{Q}$ is amenable to enforcement via projected gradient descent methods. To ensure that the orthogonal matrix $\bm{R}$ is initialized to the identity, we begin with $\bm{Q}$ as the zero matrix. The COFT approach effectively imposes two explicit constraints: one is the minimal hyperspherical energy difference, and the other directly bounds divergence from the pretrained parameters. Although traditional finetuning typically incorporates this second type of constraint implicitly, COFT treats it as an explicit requirement. While OFT achieves impressive results, COFT demonstrates enhanced finetuning stability precisely due to the explicit handling of model deviation. Figure~\ref{fig:eps_coft} illustrates the impact of varying $\epsilon$ on COFT’s performance, revealing that the parameter modulates how much the finetuned model can diverge from the pretrained version. Higher values of $\epsilon$ make the COFT-tuned model more closely match OFT, whereas reducing $\epsilon$ causes the model to maintain greater similarity to the original text-to-image diffusion baseline.

\subsection{Re-scaled Orthogonal Finetuning}

We introduce a straightforward modification to OFT whereby each neuron is endowed with an additional trainable scaling factor. This adjustment is based on the observation that altering neuron magnitudes leaves hyperspherical energy unchanged due to subsequent normalization. Formally, the forward computation adopts the form: $\thickmuskip=2mu \medmuskip=2mu\bm{z}=(\bm{R}\bm{W}^0\bm{D})^\top\bm{x}$\footnote{\textbf{Errata}: The NeurIPS camera ready submission contains an error: the re-scaled OFT’s forward pass is incorrectly listed as $\thickmuskip=2mu \medmuskip=2mu\bm{z}=(\bm{D}\bm{R}\bm{W}^0)^\top\bm{x}$. The published code used the correct computation, and all results in Appendix~\ref{app:rescaled} remain valid.}, where $\thickmuskip=2mu \medmuskip=2mu \bm{D}=\text{diag}(s_1,\cdots,s_n)\in\mathbb{R}^{n\times n}$ is a diagonal matrix with positive trainable entries $s_1,\cdots,s_n$. Unlike the original OFT described by Eq.~\eqref{eq:oft_general}, which only treats $\bm{R}$ as a parameter to be optimized, this variant also allows learning $\bm{D}$ alongside $\bm{R}$. This re-scaled extension endows OFT with greater adaptability for a negligible increase in parameter count. In our main experimental analyses, we retain the standard OFT to highlight the utility of orthogonal transformations alone; nonetheless, our findings indicate that the re-scaled version consistently performs even better (refer to Appendix~\ref{app:rescaled}).

\section{Intriguing Insights and Discussions}

\textbf{OFT is architecture-independent}. OFT is broadly applicable across various neural network structures. For example, in Transformer architectures, LoRA is typically used for adjusting the attention mechanisms~\cite{hulora2022}; therefore, for a balanced comparison with LoRA, our experiments restrict OFT’s application to attention weights. Moreover, OFT can be employed not only with fully connected layers but is also highly suitable for convolutional layers. Its block-diagonal matrix structure $\bm{R}$ brings novel interpretability to convolutions, in contrast to LoRA. When configuring one block per input channel, each block uniquely modifies a corresponding neuron channel, which conceptually parallels learning separate depth-wise convolution filters~\cite{chollet2017xception}. Alternatively, sharing all blocks in $\bm{R}$ applies the same orthogonal transformation across all channels. We offer an initial exploration into finetuning convolution layers using OFT in Appendix~\ref{app:convolution}.

\textbf{Connection to LoRA}. LoRA achieves parameter-efficient tuning by incorporating a low-rank matrix, thereby limiting drastic alterations to the pretrained weights. In contrast, OFT enforces orthogonality (full-rank) on the transformation applied to the original weight matrix, ensuring the integrity of pretraining information is preserved. The forward computation in OFT can be expressed as $\thickmuskip=2mu \medmuskip=2mu \bm{z}=(\bm{R}\bm{W}^0)^\top\bm{x}=(\bm{W}^0+(\bm{R}-\bm{I})\bm{W}^0)^\top\bm{x}$, where the term $\thickmuskip=2mu \medmuskip=2mu (\bm{R}-\bm{I})\bm{W}^0$ serves a similar role to the low-rank adjustment in LoRA. Given that $\bm{W}^0$ is generally full-rank, OFT also amounts to a low-rank adaptation when $\thickmuskip=2mu \medmuskip=2mu \bm{R}-\bm{I}$ exhibits low-rank characteristics. Analogous to LoRA's rank hyperparameter $r'$, OFT introduces a block-diagonal size $r$ to curtail the quantity of trainable weights. Notably, LoRA and OFT employ two fundamentally different approaches to achieve parameter efficiency: LoRA leverages a low-rank approximation to compress trainable parameters, while OFT leverages structural sparsity via block-diagonal orthogonal transformations.

\textbf{Why OFT converges faster}. Firstly, as demonstrated in Figure~\ref{fig:toy_ae}, semantic modification is most efficiently achieved by adjusting neuron orientations—precisely the operation enabled by OFT. Secondly, the OFT approach may be interpreted as tuning neurons constrained to a smooth hypersphere manifold, promoting a more favorable optimization surface. This theoretical advantage is corroborated by empirical findings reported in \cite{liu2021orthogonal}.

\textbf{Why not minimize hyperspherical energy}. A fundamental distinction from \cite{liu2021orthogonal} lies in our objective: we do not pursue the minimization of hyperspherical energy. In classification contexts, it is preferable for neurons to avoid redundancy. Achieving minimal hyperspherical energy would result in uniformly spaced neurons across the hypersphere, but this may be detrimental during finetuning, as it risks overwriting useful information obtained during pretraining.

\textbf{Balancing flexibility and regularity during finetuning}. Our analysis reveals an inherent compromise between the flexibility and regularity of finetuning approaches. Conventional finetuning offers maximal flexibility but suffers from instability issues and is prone to model collapse. Remarkably, OFT achieves a favorable compromise by maintaining pairwise neuron angles, despite its conceptual simplicity. The block-diagonal parameterization acts as a form of more stringent regularization imposed on the orthogonal matrix.

\textbf{No inference-time computational burden}. In contrast to ControlNet, our OFT method does not increase inference-time computational costs for the adapted model. During inference, one simply premultiplies the pretrained weight matrix $\bm{W}^0$ by the learned orthogonal transformation $\bm{R}$, forming a new equivalent parameter matrix $\thickmuskip=2mu \medmuskip=2mu \bm{W}=\bm{R}\bm{W}^0$. Consequently, inference is as efficient as with the original pretrained network.

\section{Experiments and Results}

\textbf{Experimental setup}. Experiments utilize Stable Diffusion v1.5~\cite{rombach2022high} as the foundational text-to-image model. For equitable comparison, image samples are drawn at random from each competing method. In tasks of subject-driven image synthesis, we largely replicate DreamBooth’s~\cite{ruiz2023dreambooth} methodology. In controllable synthesis experiments, protocols closely follow those described by ControlNet~\cite{zhang2023adding} and T2I-Adapter~\cite{mou2023t2i}. For a consistent comparison with LoRA, OFT or COFT is restricted to the same architectural layers where LoRA is deployed. Additional implementation details and further results can be found in Appendix~\ref{app:settings}.

\subsection{Subject-driven Generation}

\textbf{Configuration}. DreamBooth~\cite{ruiz2023dreambooth} and LoRA~\cite{hulora2022} are employed as baseline methods. All competing techniques utilize DreamBooth’s loss function. The hyperparameters for DreamBooth and LoRA are adopted from their original publications, using configurations that demonstrated optimal performance. Extended experimental results and supplementary qualitative examples are included in Appendix~\ref{app:settings},\ref{app:coft_oft},\ref{app:more_qualitative},\ref{app:failure}.

\textbf{Finetuning stability and convergence}. To begin, we assess the stability during finetuning and the convergence rates for DreamBooth, LoRA, OFT, and COFT. The corresponding results are illustrated in Figure~\ref{fig:teaser} and Figure~\ref{fig:db_convergence}. It is evident that COFT and OFT provide stable finetuning dynamics for the diffusion model. At 400 iterations, DreamBooth and OFT variants successfully enable effective control, whereas LoRA struggles to maintain subject identity. Upon reaching 2000 iterations, DreamBooth suffers from image collapse, and LoRA is unable to correctly generate a yellow shirt (instead producing yellow fur). In comparison, both OFT and COFT continue to exhibit robust and steady control over generated outputs. These findings confirm that our OFT approach ensures not only rapid but also stable convergence for subject-driven generation tasks. Furthermore, the low sensitivity of OFT and COFT to the precise finetuning iteration count significantly reduces the challenge of selecting the optimal iteration number for varying subjects. Thus, a relatively large number of iterations can be applied directly for OFT and COFT without elaborate adjustment. When $\epsilon$ is appropriately chosen for COFT, setting both the learning rate and the number of iterations becomes largely straightforward.

\textbf{Quantitative comparison}. Adopting the evaluation protocol from \cite{ruiz2023dreambooth}, we perform a quantitative analysis using metrics for subject fidelity (DINO~\cite{caron2021emerging}, CLIP-I~\cite{radford2021learning}), prompt fidelity (CLIP-T~\cite{radford2021learning}), and image diversity (LPIPS~\cite{zhang2018unreasonable}). CLIP-I involves computing the mean pairwise cosine similarity of CLIP embeddings between synthesized and real samples. DINO follows the same calculation, except employing ViT S/16 DINO embeddings in place of CLIP. The CLIP-T score captures the mean cosine similarity between CLIP embeddings of the text prompts and their corresponding generated images. For diversity, we compute the mean LPIPS cosine similarity across images generated for the same subject and prompt. According to Table~\ref{table:dreambooth_final}, both COFT and OFT surpass DreamBooth and LoRA by substantial margins in terms of DINO and CLIP-I, while delivering equivalent or superior outcomes in prompt fidelity and diversity measures. All methods are benchmarked by repeating finetuning for the same pretrained model using 30 distinct random seeds to reduce stochastic variation. Collectively, these experiments demonstrate that our OFT framework not only fosters enhanced stability and convergence but also produces superior overall performance in a consistent manner.

\textbf{Qualitative comparison}. To provide clearer insights into the advantages of OFT, we present a selection of randomly sampled outputs for subject-driven generation in Figure~\ref{fig:dreambooth_final}. For consistency, all results are produced using the same finetuned model with each method, ensuring that no prompts are individually optimized to favor a particular outcome. For every approach, we utilize the checkpoint that yields the highest validation CLIP scores. Inspection of Figure~\ref{fig:dreambooth_final} reveals that both OFT and COFT maintain strong semantic fidelity to the input subject, whereas LoRA frequently does not preserve subject identity (\emph{e.g.}, LoRA entirely loses the defining characteristics of the subject in the bowl case). In addition, OFT and COFT demonstrate more precise prompt-based control, contrasting with DreamBooth, which—despite retaining subject identity—often fails to produce images aligned with the given prompt (\emph{e.g.}, in the bowl example, first row). This qualitative evaluation highlights OFT's ability to simultaneously achieve improved prompt controllability and reliable subject preservation. Furthermore, the effectiveness of OFT remains robust to the choice of iteration count, allowing it to perform consistently even with many iterations, whereas DreamBooth and LoRA struggle under such conditions. Additional qualitative samples can be found in Appendix~\ref{app:more_qualitative}. A user study, detailed in Appendix~\ref{app:human}, provides further empirical evidence supporting OFT's superiority.

\subsection{Controllable Generation}

\textbf{Settings}. ControlNet~\cite{zhang2023adding}, T2I-Adapter~\cite{mou2023t2i}, and LoRA~\cite{hulora2022} serve as baseline comparisons. We evaluate three demanding controllable generation tasks: Canny edge to image~(C2I) using the COCO dataset~\cite{lin2014microsoft}, segmentation map to image~(S2I) on ADE20K~\cite{zhou2017scene}, and landmark to face~(L2F) on CelebA-HQ~\cite{karras2018progressive,xia2021tedigan}. Each method fine-tunes the Stable Diffusion~(SD) v1.5 model over 20 epochs for each dataset. Further results are provided in Appendix~\ref{app:more_qualitative},~\ref{app:more_control_task},~\ref{app:failure}.

\textbf{Convergence}. We assess the rate of convergence for ControlNet, T2I-Adapter, LoRA, and COFT in the L2F task, using both quantitative measurements and qualitative comparisons. For quantitative analysis, we determine the mean $\ell_2$ distance between the intended control face landmarks and the corresponding predicted landmarks. Figure~\ref{fig:face_convergence} depicts the progression of landmark errors as each model is finetuned over increasing epochs. The plots indicate that COFT and OFT both exhibit markedly accelerated convergence compared to alternatives. For example, LoRA requires 20 epochs to reach the level of performance that OFT attains by epoch 8. Notably, OFT and COFT possess a trainable parameter count similar to LoRA—which is substantially lower than that of ControlNet—yet still outperform existing techniques in convergence speed. Additional support for OFT’s rapid convergence is presented in Figure~\ref{fig:teaser}: in the rightmost example, OFT demonstrates far greater data efficiency relative to both ControlNet and LoRA, achieving satisfactory convergence using just 5\% of the ADE20K dataset. For the qualitative comparison, our analysis emphasizes OFT, COFT, and ControlNet, since the latter yields landmark error values closest to our approaches. As illustrated in Figure~\ref{fig:face_convergence_qual}, both OFT and COFT achieve stable convergence, with generated facial poses progressively matching the control landmarks. Conversely, ControlNet’s controllability appears abruptly around the 8th epoch, consistent with the sudden convergence pattern noted by \cite{zhang2023adding}. The steady convergence behavior of OFT leads to a more practical training process; the training trajectory for OFT is smoother and more predictable, thereby simplifying the search for optimal hyperparameters.

\textbf{Quantitative comparison}. We propose a metric for assessing control consistency to quantitatively measure controllable generation performance. This metric involves extracting the control signal from the output image and evaluating its correspondence with the input control signal. Specifically, for the C2I task, Intersection over Union (IoU) and the F1 score are used. In the S2I setting, we report mean IoU as well as mean and overall accuracy. The L2F task utilizes the average $\ell_2$ distance between the control and predicted landmark positions. Additional information on these consistency metrics is detailed in Appendix~\ref{app:settings}. Each method under comparison is evaluated using optimally chosen hyperparameters. Table~\ref{table:control_gen} presents the results, demonstrating that both OFT and COFT substantially outperform other approaches in terms of both strength and precision of control. Notably, adapter-based methods (\emph{e.g.}, T2I-Adapter and ControlNet) are observed to have slower convergence and inferior final performance. LoRA surpasses ControlNet on the S2I task but lags behind for C2I and L2F tasks. Overall, our analysis indicates that LoRA exhibits a comparatively lower performance ceiling, despite extensive hyperparameter optimization. Conversely, our OFT framework has yet to reach its maximal potential, as increased trainable parameters continue to improve results. We highlight that this rigorous quantitative framework for evaluating controllable generation represents a key contribution, allowing precise measurement of finetuned model control capabilities, contingent on the underlying segmentation or detection model’s accuracy.

\textbf{Qualitative comparison}. We further present a qualitative analysis of OFT and COFT in relation to leading methods such as ControlNet, T2I-Adapter, and LoRA. The randomly sampled results depicted in Figure~\ref{fig:control_final} illustrate that both OFT and COFT generate images with superior fidelity and realism while maintaining effective controllability. For the S2I task, LoRA is unable to produce outputs that adhere to the input segmentation map, whereas ControlNet, OFT, and COFT display robust control over the output. Relative to ControlNet, OFT and COFT are distinguished by their ability to synthesize visually rich and geometrically coherent images, despite employing significantly fewer model parameters. Regarding the C2I task, both OFT and COFT excel at producing plausible, semantically consistent images from rudimentary Canny edges, a scenario where T2I-Adapter and LoRA lag behind. On the L2F task, our approach offers the highest accuracy in facial pose control, performing reliably even under demanding pose conditions. Across all three control tasks, the qualitative results indicate that OFT and COFT outperform established baselines, underscoring the advantages of our OFT framework for controllable image generation. For a broader set of visual comparisons, we include additional qualitative examples across the three control tasks in Appendix~\ref{app:control_exp}, and further demonstrate the versatility of OFT on various other tasks—such as dense pose to human body, sketch to image, and depth to image—in Appendix~\ref{app:more_control_task}.

\section{Concluding Remarks and Open Problems}

Prompted by the insight that angular relationships among neurons play a vital role in visual semantic understanding, we introduce a straightforward yet powerful finetuning strategy—orthogonal finetuning—to exert control over text-to-image diffusion models. Our approach addresses two principal use cases: subject-driven generation and controllable generation. Relative to current techniques, OFT offers enhanced controllability and finetuning stability, all while requiring fewer parameters during finetuning. Crucially, OFT avoids any increase in inference cost, facilitating efficient and practical deployment.

The development of OFT raises several intriguing unresolved issues. Firstly, the orthogonality constraint in OFT is enforced through Cayley parametrization, which necessitates a matrix inversion. This requirement poses a minor barrier to the scalability of OFT. Although block diagonal parametrization helps mitigate this issue, devising a differentiable and more efficient matrix inversion method remains an open question. Secondly, OFT is uniquely positioned for compositionality: the orthogonal matrices produced by separate OFT finetuning processes can be multiplied together, and their product remains orthogonal. It remains to be explored whether this family of orthogonal matrices consistently preserves information acquired from multiple downstream tasks. Lastly, OFT's parameter efficiency is largely influenced by the chosen block diagonal structure, which can inadvertently impose additional biases and reduce adaptability. Finding ways to enhance parameter efficiency without introducing such biases represents an important avenue for future work.

\end{document}